\documentclass[10pt]{article}
\usepackage[margin=0.85in]{geometry}
\usepackage{amsmath,amssymb}
\usepackage{booktabs}
\usepackage{enumitem}
\usepackage{fancyhdr}
\usepackage{hyperref}
\usepackage{graphicx}
\usepackage{xcolor}
\usepackage{parskip}

\pagestyle{fancy}
\fancyhf{}
\lhead{ML for AIE -- Exercise Sheet 03}
\rhead{\thepage}
\renewcommand{\headrulewidth}{0.4pt}
\setlist[enumerate]{itemsep=0.2em,topsep=0.2em}
\setlist[itemize]{itemsep=0.15em,topsep=0.2em}

\newcommand{\plotplaceholder}[1]{%
  \begin{center}
    \fbox{\parbox[c][1.45in][c]{0.88\linewidth}{\centering
      \textbf{Plot placeholder}\\[0.4em]
      #1\\[0.4em]
      \emph{Replace this box with the corresponding generated plot.}}}
  \end{center}
}

\title{Exercise Sheet 03\\PCA and Model Selection: Predicting House Prices}
\author{}
\date{Winter Term 2026/27}

\begin{document}
\maketitle
\textbf{Team members:} \emph{[Name and matriculation number]} \qquad
\emph{[Name and matriculation number]}

\section{Principal Component Analysis}

\subsection{Task 1.1: Pen-and-paper calculation}

The four examples occur in opposite pairs, so their mean is zero. Their covariance matrix, using the assignment's $1/N$ convention, is
\[
\mathbf C=\frac14\sum_{i=1}^4\mathbf x^{(i)}\mathbf x^{(i)T}
=\begin{pmatrix}2.5&2\\2&2.5\end{pmatrix}.
\]
The characteristic equation is
\[
\det(\mathbf C-\lambda\mathbf I)=(2.5-\lambda)^2-4
=\lambda^2-5\lambda+2.25=0,
\]
which gives $\lambda_1=4.5$ and $\lambda_2=0.5$. Corresponding unit directions are
\[
\mathbf b_1=\frac{1}{\sqrt2}(1,1)^T,\qquad
\mathbf b_2=\frac{1}{\sqrt2}(1,-1)^T.
\]
The explained-variance ratio of $\mathbf b_1$ is $4.5/(4.5+0.5)=0.9$.

For $K=1$, the scores for the four examples are
$(3/\sqrt2,-3/\sqrt2,3/\sqrt2,-3/\sqrt2)$. Their reconstructions are respectively
$(1.5,1.5),(-1.5,-1.5),(1.5,1.5),(-1.5,-1.5)$.
Each squared Euclidean reconstruction error is $0.5$, so the mean is $0.5$,
equal to the omitted eigenvalue $\lambda_2$.

\subsection{Task 1.2: Ames Housing results and questions}

The raw feature variances range from $0.3$ (bathrooms) to $583330.6$
(lot area). Standardizing prevents a feature's units and numerical scale
from dominating the PCA merely because its variance is large. PCA is fit
to the features only; prices are used only for the score plot.

\begin{center}
\small
\begin{tabular}{rrrrr}
\toprule
$K$ & Eigenvalue $\lambda_K$ & Ratio & Cumulative ratio & Reconstruction MSE\\
\midrule
1 & 3.677 & 0.408 & 0.408 & 5.324\\
2 & 1.391 & 0.155 & 0.563 & 3.933\\
3 & 1.068 & 0.119 & 0.682 & 2.866\\
4 & 0.773 & 0.086 & 0.767 & 2.093\\
5 & 0.603 & 0.067 & 0.834 & 1.490\\
6 & 0.534 & 0.059 & 0.894 & 0.956\\
7 & 0.409 & 0.045 & 0.939 & 0.548\\
8 & 0.359 & 0.040 & 0.979 & 0.189\\
9 & 0.189 & 0.021 & 1.000 & 0.000\\
\bottomrule
\end{tabular}
\end{center}

\begin{enumerate}[label=\textbf{Q1.\arabic*:}]
\item Yes. Without standardization, lot area could dominate because its raw
variance is much larger than the other features' variances.
\item The reconstruction error using the first $K$ directions is approximately
the sum of the omitted eigenvalues. For standardized inputs, the reported
error uses $1/N$ while scikit-learn's eigenvalues use $1/(N-1)$, so the error
is approximately $(N-1)/N$ times that sum; this explains the small differences
in the table. For compression, one can retain the smallest $K$ meeting a
variance target: here $K=7$ retains $93.9\%$ of variance, while $K=6$
retains $89.4\%$.
\item All PC1 loadings are positive. The largest are living area ($0.42$),
overall quality ($0.41$), bathrooms ($0.39$), garage area ($0.37$), and
basement area ($0.34$). PC1 therefore summarizes a general combination of
house size, quality, and amenities; a larger $z_1$ indicates a larger,
better-equipped house.
\item Although PCA does not use prices, these housing characteristics are
associated with price: larger, newer, higher-quality houses tend to sell for
more. Thus the feature-based score can correlate with price, without implying
that PCA learned or causes that relationship.
\end{enumerate}

\plotplaceholder{PCA explained-variance and reconstruction-error plot
(\texttt{1\_explained\_variance.pdf})}
\plotplaceholder{First two principal-component scores colored by price
(\texttt{1\_pc\_scatter.pdf})}

\section{Overfitting and Model Evaluation}

\subsection{Task 2.1: Toy example}

For degree $1$, the fitted parameters are $\theta=(0.5,0.6)$, the training
MSE is $0.050$, and the squared error on the fifth house is $0.010$.
For degree $3$, $\theta=(-3,6.167,-2.5,0.333)$, the training MSE is
approximately zero, and the test-house squared error is $12.960$.

\begin{enumerate}[label=\textbf{Q2.\arabic*:},start=1]
\item Four distinct input values give a $4\times4$ Vandermonde design matrix
for a cubic. Its columns are linearly independent, so it has rank four and
can interpolate the four training targets exactly.
\item Training loss favors the cubic, while the held-out house favors the
line. This is overfitting: the more flexible model fits the training data
but generalizes poorly.
\end{enumerate}

\subsection{Task 2.2: Training, test, and holdout validation}

\begin{center}
\small
\begin{tabular}{rrrr}
\toprule
Degree $d$ & Training MSE & Test MSE & Holdout validation MSE\\
\midrule
1 & 780.4 & 1162.9 & 893.9\\
2 & 657.9 & 904.8 & 1210.0\\
3 & 616.6 & 903.6 & 799.2\\
4 & 585.0 & 1962.3 & 2770.7\\
5 & 582.3 & 4214.3 & 3403.9\\
6 & 574.1 & 14568.3 & 2562.0\\
\bottomrule
\end{tabular}
\end{center}
The holdout method selects degree $3$; after fitting it on all 100 training
houses, its test MSE is $903.6$.

\begin{enumerate}[label=\textbf{Q2.\arabic*:},start=3]
\item Polynomial model classes are nested: a degree-$d+1$ polynomial can
represent any degree-$d$ polynomial by setting the new coefficient to zero.
Consequently, the minimum training error cannot increase with degree.
Degree $1$ underfits relative to the better-performing degrees $2$ and $3$.
Degrees $4$--$6$ overfit: their training errors continue to decrease while
their test errors rise sharply. The high-degree fit can oscillate, particularly
near the extremes of the score range, causing large errors on unseen houses.
\item The test set must not be used to select the degree because selection
adapts the model to test outcomes and makes the reported test error optimistically
biased. The test MSE of the selected model is no longer an independent,
unbiased final estimate.
\item A validation set of 30 houses gives a noisy estimate, so a different
random split may select another degree. The selected degree should therefore
not be treated as certain; cross-validation can reduce dependence on one
particular split.
\end{enumerate}

\plotplaceholder{Polynomial fits for degrees 1, 2, and 6
(\texttt{2\_polynomial\_fits.pdf})}
\plotplaceholder{Training, holdout-validation, and test MSE by degree
(\texttt{2\_errors.pdf})}

\section{Model Selection with Cross-Validation}

The five-fold validation results and their mean are:
\begin{center}
\small
\begin{tabular}{rrrrrrr}
\toprule
$d$ & Fold 1 & Fold 2 & Fold 3 & Fold 4 & Fold 5 & CV MSE\\
\midrule
1 & 847.8 & 1329.5 & 875.8 & 487.6 & 563.4 & 820.8\\
2 & 920.6 & 1084.6 & 686.0 & 422.0 & 519.6 & 726.6\\
3 & 1141.1 & 969.1 & 646.2 & 421.4 & 451.1 & 725.8\\
4 & 1607.7 & 871.4 & 620.2 & 405.1 & 372.6 & 775.4\\
5 & 2115.1 & 952.5 & 613.7 & 415.0 & 365.6 & 892.4\\
6 & 2098.5 & 971.8 & 600.6 & 417.4 & 354.5 & 888.6\\
\bottomrule
\end{tabular}
\end{center}
Cross-validation selects $d=3$ (CV MSE $725.8$). Retraining on all
100 training houses gives test MSE $903.6$.

\begin{enumerate}[label=\textbf{Q3.\arabic*:}]
\item If each fold is used as the validation set separately, the selected
degrees are: fold 1: $d=1$; fold 2: $d=4$; fold 3: $d=6$; fold 4: $d=5$;
fold 5: $d=6$. These vary substantially, whereas averaging across folds
selects $d=3$.
\item The holdout method and cross-validation both select $d=3$ in this run.
Cross-validation is generally more trustworthy here because it averages over
five validation splits instead of relying on one 30-house set. The mean CV
errors for degrees 2 and 3 are close, however, so the selection is not
certain.
\item Using identical folds makes the degree comparison paired: each degree
is evaluated on the same held-out examples. This reduces variation due to
different splits and makes differences in validation error more attributable
to the degree.
\item Holdout is preferable when the dataset is large enough for a
representative validation set and model training is costly. It requires fewer
fits than $K$-fold CV, while CV is useful when data is limited and a more
stable estimate is worth the additional computation.
\end{enumerate}

\plotplaceholder{Cross-validation and test MSE by degree; selected degree
(\texttt{3\_cross\_validation.pdf})}

\section*{Generative AI disclosure}
An AI coding assistant was used to help implement and check the analysis and
to draft this report. The reported numerical results were checked against
the provided Python task outputs. AI-generated code and explanations can
still contain errors; the authors are responsible for understanding and
verifying the submitted solution.

\end{document}
